\documentclass[11pt]{article}
\usepackage[margin=1in]{geometry}
\usepackage{amsmath,amssymb,amsthm}
\usepackage{booktabs}
\usepackage{array}
\usepackage{hyperref}
\usepackage{microtype}

\newtheorem{theorem}{Theorem}
\newtheorem{proposition}{Proposition}
\newtheorem{corollary}{Corollary}
\newtheorem{definition}{Definition}
\theoremstyle{remark}
\newtheorem*{remark}{Remark}

\title{Beaufort Keystream Anomalies in Kryptos K4}
\author{Colin Patrick \and Claude (Anthropic)}
\date{March 31, 2026 (revised)}

\begin{document}
\maketitle

\begin{abstract}
We report model-independent statistical anomalies in the Beaufort
keystream of Kryptos K4 at its 24~known crib positions.
The keystream values $k[i] = (C[i] + P[i]) \bmod 26$ are enriched
for the 7-letter set $\Pi = \{B,G,I,K,O,W,Z\}$: 13/24 overall
($p = 4.3 \times 10^{-3}$), with extreme concentration at
BERLINCLOCK positions 63--70 (7/8, $p = 6.3 \times 10^{-4}$).
These letters correspond to columns~0 and~3 of the Kryptos Alphabet
arranged as a $5$-wide Polybius grid.
Bean~(2021) independently found that 13/24 keystream values are
multiples of~5 in a reversed-KA indexing ($p \approx 9.5 \times 10^{-4}$),
sharing a mod-5 connection with the grid.
All keystream findings depend only on the carved ciphertext and the
known cribs---no cipher model, keyword, or null mask is assumed.
A previously reported null-palette restriction (7~distinct letters at
17~consensus null positions) has been \textbf{downgraded}: cross-model
testing shows the restriction is an artifact of the iterative
optimization procedure, not a model-invariant property of K4
(Section~\ref{sec:downgrade}).
\end{abstract}

%% ================================================================
\section{Setup}
%% ================================================================

Let $\mathcal{A} = \{A, B, \ldots, Z\}$ with standard indexing $A{=}0,
\ldots, Z{=}25$.  All arithmetic modulo~26.

\begin{definition}[Kryptos Alphabet]
$\mathrm{KA} = {}${\small\texttt{KRYPTOSABCDEFGHIJLMNQUVWXZ}}, with
$\mathrm{KA}[0]{=}K$, $\mathrm{KA}[1]{=}R$, \ldots, $\mathrm{KA}[25]{=}Z$.
All 26~letters appear exactly once.
\end{definition}

\begin{definition}[K4 Ciphertext and Cribs]
{\small
\begin{align*}
C &= \texttt{OBKRUOXOGHULBSOLIFBBWFLRVQQPRNGKSSOTWTQSJQSSEKZZWAT} \\
  &\phantom{{}={}}\texttt{JKLUDIAWINFBNYPVTTMZFPKWGDKZXTJCDIGKUHUAUEKCAR}
\end{align*}
}
with $|C| = 97$ and 0-indexed positions.  Known cribs (24~characters total):
\begin{align*}
P[21..33] &= \texttt{EASTNORTHEAST} \quad (13\text{ chars}), \\
P[63..73] &= \texttt{BERLINCLOCK} \quad\;\; (11\text{ chars}).
\end{align*}
\end{definition}

\begin{definition}[$5$-Wide KA Polybius Grid]\label{def:grid}
Arrange KA into a grid with 5~columns:
\[
\begin{array}{c|ccccc}
      & c_0 & c_1 & c_2 & c_3 & c_4 \\ \hline
r_0 & \mathbf{K} & R & Y & P & T \\
r_1 & \mathbf{O} & S & A & \mathbf{B} & C \\
r_2 & D & E & F & \mathbf{G} & H \\
r_3 & \mathbf{I} & J & L & M & N \\
r_4 & Q & U & V & \mathbf{W} & X \\
r_5 & \mathbf{Z} &   &   &   &
\end{array}
\]
where bold letters occupy columns~0 and~3 exclusively.
For letter $\ell$, define $\mathrm{col}(\ell) = \mathrm{KA\text{-}idx}(\ell) \bmod 5$.
\end{definition}

\begin{definition}[Palette]
$\Pi = \{\ell \in \mathcal{A} : \mathrm{col}(\ell) \in \{0, 3\}\}
     = \{B, G, I, K, O, W, Z\}$.
\end{definition}

%% ================================================================
\section{Proposition 1: Beaufort Keystream Enrichment}
%% ================================================================

\begin{definition}[Beaufort Keystream at Crib Positions]
For each known crib position~$i$:
\[
k[i] = (C[i] + P[i]) \bmod 26.
\]
This value is \textbf{model-independent}: it depends only on the
carved text and the known plaintext, not on any cipher model,
keyword, or null mask.
\end{definition}

\begin{proposition}[Palette Enrichment]\label{prop:keystream}
Of the 24 Beaufort keystream values at crib positions, 13 are
members of~$\Pi$:
\[
|\{i \in \mathcal{C} : k[i] \in \Pi\}| = 13.
\]
\end{proposition}

\begin{proof}
Direct computation at all 24~crib positions.

\medskip
\noindent\textit{EASTNORTHEAST} (positions 21--33):
\begin{center}
\begin{tabular}{ccccl}
\toprule
Pos & $C[i]$ & $P[i]$ & $k[i]$ & $\in \Pi$? \\
\midrule
21 & F & E & J\,(9)  & No \\
22 & L & A & L\,(11) & No \\
23 & R & S & J\,(9)  & No \\
24 & V & T & O\,(14) & Yes \\
25 & Q & N & D\,(3)  & No \\
26 & Q & O & E\,(4)  & No \\
27 & P & R & G\,(6)  & Yes \\
28 & R & T & K\,(10) & Yes \\
29 & N & H & U\,(20) & No \\
30 & G & E & K\,(10) & Yes \\
31 & K & A & K\,(10) & Yes \\
32 & S & S & K\,(10) & Yes \\
33 & S & T & L\,(11) & No \\
\bottomrule
\end{tabular}
\end{center}
ENE palette hits: $6/13$.

\medskip
\noindent\textit{BERLINCLOCK} (positions 63--73):
\begin{center}
\begin{tabular}{ccccl}
\toprule
Pos & $C[i]$ & $P[i]$ & $k[i]$ & $\in \Pi$? \\
\midrule
63 & N & B & O\,(14) & Yes \\
64 & Y & E & C\,(2)  & No \\
65 & P & R & G\,(6)  & Yes \\
66 & V & L & G\,(6)  & Yes \\
67 & T & I & B\,(1)  & Yes \\
68 & T & N & G\,(6)  & Yes \\
69 & M & C & O\,(14) & Yes \\
70 & Z & L & K\,(10) & Yes \\
71 & F & O & T\,(19) & No \\
72 & P & C & R\,(17) & No \\
73 & K & K & U\,(20) & No \\
\bottomrule
\end{tabular}
\end{center}
BCL palette hits: $7/11$.  Total: $13/24$.
\end{proof}

\paragraph{Statistical test.}
Under $H_0$ (keystream uniform over $\mathcal{A}$),
$P(\text{palette letter}) = |\Pi|/26 = 7/26 \approx 0.269$.
\[
P(X \geq 13 \mid n{=}24,\; p{=}7/26) = 4.3 \times 10^{-3}.
\]
For the BCL first 8 positions specifically:
\[
P(X \geq 7 \mid n{=}8,\; p{=}7/26) = 6.3 \times 10^{-4}.
\]

\paragraph{Variant specificity.}
Only Beaufort ($A{=}0$) achieves 13/24.  Vigen\`ere gives 9/24;
Variant Beaufort gives 5/24.

\paragraph{Search burden.}
Three cipher variants were tested.  After Bonferroni correction
($\times 3$), the BCL first-8 result gives $p = 1.9 \times 10^{-3}$,
still significant.  The palette $\Pi$ was identified from the
null-position analysis (Section~\ref{sec:downgrade}), making the
keystream test partially pre-specified (the palette was defined
before the keystream was examined, on disjoint positions) but
ultimately derived from the same ciphertext.

%% ================================================================
\section{Bean's Mod-5 Observation (2021)}\label{sec:bean}
%% ================================================================

\begin{proposition}[Bean's Mod-5 Observation]\label{prop:bean}
Define the reversed Kryptos Alphabet
$\mathrm{rKA} = {}${\small\texttt{ZXWVUQNMLJIHGFEDCBASOTPYRK}} and compute:
\[
\beta[i] = (\mathrm{rKA\text{-}idx}(C[i]) - \mathrm{idx}(P[i])) \bmod 26
\]
at each crib position~$i$.  Then $13$ of $24$ values satisfy
$\beta[i] \equiv 0 \pmod{5}$.
\end{proposition}

\begin{proof}
Bean (2021, footnote~3).  We independently verify:
$P(X \geq 13 \mid n{=}24,\; p{=}6/26) = 9.5 \times 10^{-4}$.
\end{proof}

\paragraph{Apparent structural connection (incorrect).}
An earlier draft of this paper claimed that
$\beta[i] \equiv 0 \pmod{5}$ is equivalent to
$\mathrm{col}(C[i]) \equiv -\mathrm{idx}(P[i]) \pmod{5}$,
which would directly select KA Polybius columns.
This equivalence is \textbf{wrong}: it requires that
$(x + y) \bmod 26 \equiv (x + y) \bmod 5$,
which holds only when $x + y < 26$.
Since $26 \equiv 1 \pmod{5}$, reducing mod~26 before mod~5
introduces a $\pm 1$ shift for values $\geq 26$.  Empirically,
only 9 of the 13 Bean mod-5 hits satisfy the column congruence.

\paragraph{Relationship to the Polybius grid.}
Bean's mod-5 and the palette enrichment share a surface
connection: both involve divisibility by~5, and the KA Polybius
grid has 5~columns.  However, the algebraic reduction from
Bean's statistic to column selection does not hold exactly.
Bean's observation is a genuine, independently significant
anomaly ($p = 9.5 \times 10^{-4}$) that is \emph{related to}
but not \emph{equivalent to} the Polybius column structure.

\paragraph{Independence.}
Bean's observation was published in 2021, before the palette
$\Pi$ was identified.  The two findings were made by different
researchers, from different starting points (reversed-KA
keystream arithmetic vs.\ null-character identity).

%% ================================================================
\section{K-Sum Bias}
%% ================================================================

\begin{proposition}[Keystream K-Sum Deficit]\label{prop:ksum}
The Beaufort keystream at the 24~crib positions is
{\small\texttt{JLJODEGKUKKKLOCGGBGOKTRU}}.
Mapping each letter to its KA index and summing:
\[
\Sigma = \sum_{i \in \mathcal{C}} \mathrm{KA\text{-}idx}(k[i]) = 218.
\]
Under a uniform keystream, $E[\Sigma] = 24 \cdot 25/2 = 300$,
$\mathrm{SD}(\Sigma) = \sqrt{24 \cdot (26^2 - 1)/12} \approx 36.7$.
The deficit is $82$, giving $z = -2.23$ ($p \approx 0.013$,
one-tailed).
\end{proposition}

\paragraph{Conditional status.}
The K-sum deficit is largely a consequence of the palette enrichment
(Proposition~\ref{prop:keystream}): palette letters $\Pi$ occupy
early KA positions (mean KA-index 7.7 vs.\ 12.5 overall), so any
keystream enriched for $\Pi$ will mechanically produce a low KA-sum.
The K-sum is therefore a \emph{conditional redescription} of the
palette signal, not an independent observation.

\paragraph{Note.}
This observation was first noted by Bean in private correspondence
(March 2026), who reported the K-sum as ``the only interesting
thing about that key.''

%% ================================================================
\section{Downgraded: Null Palette Restriction}\label{sec:downgrade}
%% ================================================================

An earlier version of this paper presented the null palette
restriction as a co-equal finding alongside the keystream enrichment.
Cross-model testing has revealed this claim to be \textbf{substantially
weaker than originally reported}.

\paragraph{Original claim.}
The 17~consensus null positions (identified by simulated annealing
under a KA autokey Vigen\`ere model) contain only 7~distinct letters,
all from $\Pi$.  Monte Carlo testing gave $p \approx 6.3 \times 10^{-5}$.

\paragraph{What went wrong.}
The 7-distinct result was produced by a two-phase optimization:
(1)~200~SA restarts identify high-frequency positions;
(2)~those positions are \emph{fixed}, and SA optimizes the remaining
positions subject to that constraint.  Phase~2 mechanically restricts
which letters can appear, because fixing early-selected positions
removes letter diversity that would otherwise accumulate.

When the full 200-restart protocol is run \emph{without} the
bootstrapping phase, even the original discovery model (KA autokey
Vigen\`ere) produces 13~distinct letters at its consensus positions
---not~7.  Across 5~cipher models tested:

\begin{center}
\begin{tabular}{lcc}
\toprule
Model & Distinct & $|\Pi \cap \text{palette}|$ \\
\midrule
KA autokey Vigen\`ere & 13 & 6/7 \\
KA autokey Beaufort & 10 & 3/7 \\
KA periodic Beaufort & 10 & 6/7 \\
AZ periodic Vigen\`ere & 13 & 4/7 \\
AZ periodic Beaufort & 14 & 4/7 \\
\bottomrule
\end{tabular}
\end{center}

Zero models produce $\leq 7$ distinct letters.  The mean position
Jaccard similarity across models is $0.152$, confirming that
positions are heavily model-dependent.

\paragraph{What survives.}
The 0/500 shuffled-CT result remains valid: when K4's letter
arrangement is randomly permuted, SA-optimized masks produce
$9$--$17$ distinct letters (mean $12.7$), never $\leq 7$.
This means real K4 is unusual in \emph{permitting} low-diversity
subsets under aggressive optimization---but the optimization
procedure itself is responsible for driving diversity below~10.

\paragraph{Implication for the joint significance.}
The previously reported joint $p = 1.4 \times 10^{-7}$ (combining
palette restriction and keystream enrichment) is \textbf{no longer
valid}, because the palette restriction component has been
invalidated.  The keystream enrichment
(Proposition~\ref{prop:keystream}, $p = 4.3 \times 10^{-3}$;
BCL first-8: $p = 6.3 \times 10^{-4}$) stands on its own as
the primary surviving signal.

%% ================================================================
\section{Summary of Evidence}
%% ================================================================

\begin{center}
\begin{tabular}{llll}
\toprule
Observation & Status & $p$-value & Model-free? \\
\midrule
Keystream 13/24 $\in \Pi$ & \textbf{Active} & $4.3 \times 10^{-3}$ & Yes \\
BCL first 8: 7/8 $\in \Pi$ & \textbf{Active} & $6.3 \times 10^{-4}$ & Yes \\
Bean mod-5 (2021) & \textbf{Active} & $9.5 \times 10^{-4}$ & Yes \\
K-sum $= 218$ & Conditional & $0.013$ & Yes (but not independent) \\
Beaufort $>$ other variants & \textbf{Active} & Bonferroni $\times 3$ & Yes \\
\midrule
Null palette $\leq 7$ distinct & \textbf{Downgraded} & was $6.3 \times 10^{-5}$ & No \\
Joint $p = 1.4 \times 10^{-7}$ & \textbf{Retracted} & --- & No \\
\bottomrule
\end{tabular}
\end{center}

The surviving evidence is moderate: three model-independent signals
(keystream palette enrichment, BCL concentration, Bean mod-5) all
point toward a mod-5 / KA-column structure in the Beaufort keystream.
None individually reaches $p < 10^{-4}$ after search-breadth
correction, and none produces a decryption.  These are descriptive
anomalies that constrain the mechanism but do not identify it.

%% ================================================================
\section{Open Questions}
%% ================================================================

\noindent\textbf{Resolved: Width-21 bigram is a stego artifact.}
The width-21 bigram signal ($p = 1.6 \times 10^{-4}$) is present in
the raw 97-character text but \emph{disappears completely} after
removing the 17~consensus null positions.

\medskip
\noindent\textbf{Resolved: Width-10/17 bigrams are cipher-layer artifacts.}
On the null-extracted text, bigram enrichment at widths~10
and~17 ($p = 6$--$8 \times 10^{-3}$) is destroyed by undoing a
column-7 transposition, suggesting mathematical artifact rather
than cipher-layer structure.

\medskip
\noindent\textbf{Open: Why is the Beaufort keystream enriched for $\Pi$?}
The enrichment is a property of the carved ciphertext at crib
positions.  Any correct theory of K4 must explain why
$(C[i] + P[i]) \bmod 26$ lands in $\Pi$ at 7/8 BERLINCLOCK
positions.  Under Beaufort, this means the operational key at
those positions is almost entirely drawn from~$\Pi$.

\medskip
\noindent\textbf{Open: Running-key source identification.}
Running key is the only structured non-periodic key model surviving
Bean constraints.  Keystream fragments at crib positions do not
resemble English under identity or columnar transposition (widths
6, 8, 9).  Running key with a monoalphabetic inner layer remains
underdetermined.

%% ================================================================
\section*{Reproducibility}
%% ================================================================

All computations use the \texttt{kryptos} kernel
({\small\texttt{src/kryptos/kernel/constants.py}}, which self-verifies CT,
cribs, and Bean constraints at import time).
Cross-model palette stability:
{\small\texttt{scripts/statistical/e\_cross\_model\_palette\_stability\_01.py}},
5 models $\times$ 200 SA restarts $\times$ 300K steps.
Running-key English-likelihood test:
{\small\texttt{scripts/statistical/e\_keystream\_english\_likelihood\_01.py}},
seed \texttt{42}, $N = 1{,}000{,}000$ null fragments, 17{,}124 columnar orderings.
Source: {\small\texttt{github.com/jcolinpatrick/kryptos}}.
As of March~31, 2026: 940~scripts tracked, 390~eliminations,
$671 \times 10^9$+ configurations tested.

\begin{thebibliography}{9}
\bibitem{bean2021} R.~Bean, ``Cryptodiagnosis of Kryptos K4,''
  \emph{Proceedings of the 4th International Conference on Historical
  Cryptology (HistoCrypt 2021)}, pp.~13--21.
\bibitem{sanborn1990} J.~Sanborn, Dedication speech, November 3, 1990.
  RR~Auction lot documentation.
\end{thebibliography}

\end{document}
